% Generated draft: requires mathematical and source review.
\documentclass[11pt]{article}
\usepackage[margin=25mm]{geometry}
\usepackage{fontspec}
\setmainfont{Latin Modern Roman}
\usepackage{amsmath,amssymb,mathtools,unicode-math}
\setmathfont{Latin Modern Math}
\usepackage{enumitem,needspace,xurl,hyperref}
\setlength{\parindent}{0pt}
\setlength{\parskip}{5pt}
\setlength{\emergencystretch}{3em}
\newcommand{\N}{\mathbb{N}}
\newcommand{\Z}{\mathbb{Z}}
\newcommand{\Q}{\mathbb{Q}}
\newcommand{\R}{\mathbb{R}}
\newcommand{\C}{\mathbb{C}}
\newcommand{\F}{\mathbb{F}}
\newcommand{\A}{\mathbb{A}}
\newcommand{\PP}{\mathbb{P}}
\newcommand{\E}{\mathbb{E}}
\newcommand{\eps}{\varepsilon}
\newcommand{\dd}{\,\mathrm{d}}
\begin{document}
\section*{UnsolvedMath ID: 30002167}
\textbf{Source problem code:} catalogue entry
\subsection{Definitions and mathematical statement}
Work in the Euclidean plane \(\R^2\), with Euclidean norm \(\|\cdot\|_2\). Write
\[
B_{\mathrm{cl}}=\{x\in\R^2:\|x\|_2\leq 1\},\qquad
B_{\mathrm{op}}=\{x\in\R^2:\|x\|_2<1\}.
\]
The supplied source calls \(B\) the unit disk without explicitly specifying whether its boundary is included. The questions below are therefore stated with \(B\) denoting the intended one of these two disks; that boundary convention remains to be confirmed.

For a finite set \(P=\{p_1,\ldots,p_n\}\subseteq B\) with \(n\geq 3\), a Hamiltonian cycle means a cyclic ordering \(p_{\sigma(1)},\ldots,p_{\sigma(n)}\) of all its points, with straight edges between consecutive points and from the last point to the first. Its squared-edge cost is
\[
H(P,\sigma)=\sum_{i=1}^{n}\bigl\|p_{\sigma(i+1)}-p_{\sigma(i)}\bigr\|_2^2,
\qquad \sigma(n+1)=\sigma(1),
\]
where \(\sigma\) is a permutation of \(\{1,\ldots,n\}\). No restriction on edge crossings is imposed in this graph-theoretic formulation; the supplied record does not explicitly discuss crossings.

\textbf{Hamiltonian-cycle question.} Is it true that for every integer \(n\geq 3\) and every \(n\)-point set \(P\subseteq B\), there exists such a permutation \(\sigma\) satisfying
\[
H(P,\sigma)\leq 8?
\]
The source says only that \(P\) is finite. Its intended cycle convention for \(|P|\in\{0,1,2\}\) is not supplied, so those cases are not included in the displayed formulation.

For a finite set \(P\) of even cardinality, a perfect matching is a collection \(M\) of pairwise disjoint two-element subsets of \(P\) whose union is \(P\). Define its squared-edge cost by
\[
W(M)=\sum_{\{p,q\}\in M}\|p-q\|_2^2.
\]
\textbf{Perfect-matching question.} Is it true that for every integer \(m\geq 0\) and every \(2m\)-point set \(P\subseteq B\), there exists a perfect matching \(M\) satisfying
\[
W(M)\leq 4?
\]
Here the empty set has the empty perfect matching, with cost zero.

\textbf{Requested extension.} Determine analogous bounds for finite point sets in other convex figures \(C\). The cleaned statement describes these as convex bodies; a conventional planar convex body is a compact convex subset of \(\R^2\) with nonempty interior. The original statement asks for analogues with \(C\) in place of the disk ``or the unit square.'' The supplied record gives no unit-square bound, precise class of convex figures, normalization of \(C\), or proposed dependence of the bounds on \(C\). Thus this extension is an open-ended request, not a uniquely specified additional inequality.
\subsection{Short English statement}
For every finite set of points in the Euclidean unit disk, ask whether a Hamiltonian cycle can have total squared edge length at most \(8\), and, when the number of points is even, whether a perfect matching can have total squared edge length at most \(4\). The cycle statement above covers at least three points. Boundary conventions, smaller cycle cases, and the intended bounds for other convex figures require clarification. This is a statement-only draft requiring human mathematical review.
\subsection{Sources}
\begin{itemize}
\item Local UnsolvedMath registry, permanent ID \(30002167\), title \emph{Short Hamiltonian Cycles and Matchings in Convex Bodies}, source problem code \texttt{catalogue entry}, in the supplied Oberwolfach Reports Open Problems collection. No individual proposer is identified; the supplied proposed year is 2012.
\item Supplied report citation: \emph{Optimal and Near Optimal Configurations on Lattices and Manifolds} (2013), Oberwolfach Reports, pp.\ 2429--2492, DOI \texttt{10.4171/owr/2012/40}. Supplied source URL and literature-source link labeled ``OWR report 40'': \url{https://doi.org/10.4171/owr/2012/40}. The record attributes the constants, parity condition, and convex-body extension to the official open-problem session; that attribution has not been independently verified here.
\item Supplied literature-source link labeled ``Current open problems survey'': \url{https://www.researchgate.net/publication/328301196_Current_Open_Problems_in_Discrete_and_Computational_Geometry}. The record supplies no authors or complete bibliographic citation for this link.
\item Supplied literature assessment dated exactly \texttt{2026-08-21}: the disk cycle and matching questions are described as remaining open in later discrete-geometry surveys, with exact constants and extensions to other convex figures unsettled. This is supplied dated curation, not an independently verified literature assessment. The background explicitly requests independent verification of time-sensitive status and cited claims.
\item The supplied registry labels the problem \texttt{open} and the statement \texttt{corrected\_verified}. These are source metadata labels, not mathematical or source verification performed for this draft.
\end{itemize}
\end{document}
